\documentclass[12pt,a4paper,notitlepage,twocolumn]{article}
\usepackage{amsmath}
\usepackage{amssymb}
\usepackage{amsbsy}
\usepackage{float}
\usepackage[french]{babel}

\usepackage{palatino}

 \usepackage[active]{srcltx}
\usepackage{scrtime}

\newcounter{qnumber}
\setcounter{qnumber}{0}

\newcounter{enumber}
\setcounter{enumber}{0}

\newcommand{\question}{\textbf{(\addtocounter{qnumber}{1}\theqnumber)}\,}
\newcommand{\exercice}{\textbf{\addtocounter{enumber}{1}\textsc{Exercice} \theenumber}.\,}
\setlength{\parindent}{0cm}


\begin{document}

\title{\textsc{Économétrie Approfondie}}
\date{Jeudi 18 décembre 2025}

\maketitle

\thispagestyle{empty}

\begin{quote}
  \textit{Les réponses non commentées ou insuffisamment détaillées ne seront pas
    considérées. Prenez le temps de faire des phrase.}
\end{quote}

\bigskip


\exercice On suppose que les données sont générées par le modèle suivant~:
\[
  y_i = x_{i,1}\beta_1 + \ldots x_{i,K}\beta_K + \varepsilon_i
\]
où $x_{i,k}$, pour $k=1,\ldots,K$ sont des variables explicatives
déterministes, $\beta_k$, pour $k=1,\ldots,K$, sont des paramètres
réels, $\varepsilon_i$ une variable aléatoire i.i.d. centrée de
variance $\sigma_{\varepsilon}^2$. \question Donner la représentation
matricielle de ce modèle. \question Définir et donner l'expression de
l'estimateur des MCO du vecteur de paramètres $\beta$, que l'on notera
$\hat b$ (en montrant comment on arrive à cette expression et en
explicitant les hypothèses nécessaires pour que cet estimateur
existe). \question Montrer que cet estimateur est sans
biais. \question Calculer la variance de cet estimateur. \question
Pourquoi l'estimateur des MCO a t-il une variance~? \question Montrer
que l'estimateur des MCO $\hat b$ convergent en probabilité vers $\beta$.\newline

\setcounter{qnumber}{0}
\bigskip

\exercice On considère le modèle générateur des données suivant~:
\[
y_i = x_i \beta + \varepsilon_i
\]
pour $i=1,\ldots,N$, avec $\beta$ un vecteur colonne de paramètres
réels, et $x_i$ un vecteur ligne de variables explicatives
déterministes. On suppose que les erreurs $\varepsilon_i$ sont normalement distribuées et que
\[
\varepsilon = (\varepsilon_1, \ldots, \varepsilon_N)^T \sim \mathcal{N}(\mathbf{0}, \Omega)
\]
où $\Omega$ est une matrice symétrique et définie positive. On considère le modèle empirique~:
\[
  y_i = x_i b + \epsilon_i
 \]
 \question L'estimateur des MCO, $\hat b$, est-il un estimateur sans
 biais~? Justifier. \question Calculer la variance de l'estimateur des
 MCO . \question L'estimateur des MCO est-il efficace parmi la classe
 des estimateurs linéaires et sans biais~? Justifier. \question
 Proposer un estimateur efficace et calculer sa variance, en supposant
 que la matrice $\Omega$ est connue. \question Que faire si la matrice
 $\Omega$ est inconnue~? Proposez un nouvel estimateur. \question
 Est-il possible d'estimer ce modèle par maximum de
 vraisemblance~?\newline

\setcounter{qnumber}{0}
\bigskip

\exercice Sur le marché d'un bien la demande et l'offre à la date $t$ sont données par~:
\begin{equation}\label{s1}
  \begin{cases}
    y_t^d &= \alpha_0 + \alpha_1 P_t + \alpha_2 R_t + u_t\\
    y_t^o &= \gamma_0 + \gamma_1 P_t + \gamma_2 W_t + v_t
  \end{cases}
\end{equation}
où $y_t^d$ et $y_t^o$ mesurent la demande et l'offre de bien, $P_t$ le
prix du bien, $R_t$ le revenu des consommateurs, $W_t$ les coûts de
production, $u_t$ et $v_t$ sont des variables aléatoires iid,
centrées, de variances $\sigma_u^2$ et $\sigma_v^2$ et non corrélées,
qui s'interprètent comme des chocs de demande et d'offre. Le système
d'équation \ref{s1} est le modèle structurel. On suppose que les
échanges ne se font qu'à l'équilibre, c'est à dire lorsque la quantité
demandée est égale à la quantité offerte. On suppose que les coûts de
production et le choc de demande sont orthogonaux. On suppose que la
variable de revenu est orthognale aux chocs. \question Quels sont
les signes attendus des paramètres~? \question Calculer le prix
d'équilibre pour montrer qu'il s'agit d'une variable
aléatoire. \question Un économètre souhaite estimer l'équation de
demande par les MCO. Pensez-vous que cela soit une bonne idée~?
Justifier. \question Proposer une méthode d'estimation permettant
d'obtenir un estimateur convergent des paramètres de l'équation de
demande. \question Cet estimateur est-il sans biais~?

\end{document}

%%% Local Variables:
%%% mode: latex
%%% TeX-master: t
%%% End:
